\subsection{测度与连续函数的乘积}

设 X 是一个局部紧空间，g 是 X 到 C 中的一个连续映射。显然 \(f \mapsto gf\) 是 \(\mathcal{K}(\mathrm{X};\mathbf{C})\) 到其自身中的一个线性映射；我们来证明这一映射是连续的。事实上，对 X 的每个紧子集 K 及每个函数 \(f \in \mathcal{K}(\mathrm{X},\mathrm{K};\mathbf{C})\) ，有 \(gf \in \mathcal{K}(\mathrm{X},\mathrm{K};\mathbf{C})\) ；此外，若 \(b_{\mathrm{K}} = \sup_{x \in \mathrm{K}} |g(x)|\) ，则 \(\|gf\| \leqslant b_{\mathrm{K}} \|f\|\) ，由此即得我们的断言（TVS, II, §4, No. 4, 命题 5）。于是这一连续线性映射的转置（TVS, II, §6, No. 4）是 \(\mathcal{M}(\mathrm{X};\mathbf{C})\) 到其自身中的一个线性映射，记作 \(\mu \mapsto g \cdot \mu\) （如不致混淆，也记作 \(\mu \mapsto g\mu\) ）。若 \(\nu = g \cdot \mu\) ，则对每个函数 \(f \in \mathcal{K}(\mathrm{X};\mathbf{C})\) ，

\[
\langle f, \nu \rangle = \langle g f, \mu \rangle\tag{6}
\]

或者写作

\[
\int f (x) d \nu (x) = \int f (x) g (x) d \mu (x)
\]

（它缩写为形式 \(d\nu(x) = g(x) \, d\mu(x)\) ）。称 \(g \cdot \mu\) 为测度 \(\mu\) 与函数 \(g\) 的乘积，或称为以 \(g\) 为密度（关于 \(\mu\) ）的测度（参见 Ch. V, §5, No. 2, 定义 2）。若 \(g_1, g_2\) 是 X 到 C 中的两个连续映射，且 \(\mu_1, \mu_2\) 是 X 上的两个测度，则

\[
\begin{array}{c} (g _ {1} + g _ {2}) \cdot \mu = g _ {1} \cdot \mu + g _ {2} \cdot \mu , \quad g \cdot (\mu_ {1} + \mu_ {2}) = g \cdot \mu_ {1} + g \cdot \mu_ {2}, \\ (g _ {1} g _ {2}) \cdot \mu = g _ {1} \cdot (g _ {2} \cdot \mu). \end{array}
\]

此外，\(1 \cdot \mu = \mu\) （这里 1 表示 X 上恒等于 1 的常值函数）；集合 \(\mathcal{M}(\mathrm{X};\mathbf{C})\) 赋以外合成律 \((g,\mu) \mapsto g \cdot \mu\) 及其加法结构后，是环 \(\mathcal{C}(\mathrm{X};\mathbf{C})\) 上的一个模。

